% =========================================================
% Chapter 5: The Laws of Motion
% POSN 1 Physics Preparation Notes
% Language: Thai
% Compiler: XeLaTeX
% =========================================================

\chapter{กฎการเคลื่อนที่}

\begin{center}
    {\Large \textbf{The Laws of Motion}}\\[4pt]
    {\normalsize เอกสารเตรียมสอบคัดเลือก สอวน. ฟิสิกส์ ค่าย 1}\\[6pt]
    {\normalsize จัดทำโดย \textbf{Tames Thanakrit Damduan}}\\[2pt]
    {\footnotesize \textcopyright\ 2026 Tames Thanakrit Damduan. All rights reserved.}\\[8pt]
\end{center}

\section*{เป้าหมายของบทเรียน}

หลังจากเรียนบทนี้ นักเรียนควรทำสิ่งต่อไปนี้ได้

\begin{enumerate}
    \item เข้าใจว่าแรงเป็นสาเหตุของการเปลี่ยนสภาพการเคลื่อนที่
    \item จำแนกแรงสัมผัสและแรงสนามได้
    \item ใช้กฎข้อที่หนึ่งของนิวตันอธิบายสมดุลและกรอบอ้างอิงเฉื่อยได้
    \item เข้าใจมวลในฐานะการวัดความเฉื่อยของวัตถุ
    \item ใช้กฎข้อที่สองของนิวตัน $\sum \vec{F}=m\vec{a}$ ได้ถูกต้อง
    \item แยกแรงเป็นองค์ประกอบตามแกน $x$ และ $y$ ได้
    \item เขียนแผนภาพวัตถุอิสระหรือ free-body diagram ได้
    \item เข้าใจแรงโน้มถ่วง น้ำหนัก แรงปกติ แรงตึงเชือก และแรงเสียดทาน
    \item ใช้กฎข้อที่สามของนิวตันแยกคู่แรงกิริยาและปฏิกิริยาได้
    \item แก้โจทย์สมดุล โจทย์วัตถุเร่ง โจทย์พื้นเอียง และโจทย์แรงเสียดทานระดับ สอวน. ได้
\end{enumerate}

\begin{tcolorbox}[title=แนวคิดสำคัญของบทนี้]
บทก่อนหน้านี้ตอบคำถามว่า ``วัตถุเคลื่อนที่อย่างไร'' แต่บทนี้ตอบคำถามว่า ``ทำไมวัตถุจึงเคลื่อนที่แบบนั้น'' หลักสำคัญคือให้เริ่มจากวาดแรงทุกแรงที่กระทำต่อวัตถุ แล้วใช้กฎของนิวตันแยกตามแกน ไม่ใช่เริ่มจากการจำสูตร
\end{tcolorbox}

\newpage{}

% =========================================================
\section{แนวคิดของแรง}

\subsection{แรงคืออะไร}

ในชีวิตประจำวัน เราเห็นแรงจากการผลัก ดึง ยก เตะ หรือกดวัตถุ แต่ในฟิสิกส์ แรงมีความหมายกว้างกว่านั้น

แรงคือปฏิสัมพันธ์ที่สามารถทำให้วัตถุเปลี่ยนสภาพการเคลื่อนที่ได้ เช่น ทำให้วัตถุเริ่มเคลื่อนที่ หยุด ช้าลง เร็วขึ้น หรือเปลี่ยนทิศทาง

\begin{tcolorbox}[title=นิยาม]
แรง คือ ปฏิสัมพันธ์ที่สามารถทำให้ความเร็วของวัตถุเปลี่ยนได้

แรงเป็นเวกเตอร์ จึงต้องมีทั้งขนาดและทิศทาง
\end{tcolorbox}

แรงใช้หน่วยนิวตัน

\[
\si{N}
\]

โดย

\[
1\,\si{N}=1\,\si{kg.m/s^2}
\]

\subsection{แรงสัมผัสและแรงสนาม}

แรงแบ่งได้เป็นสองกลุ่มใหญ่

\[
\text{แรงสัมผัส}
\]

และ

\[
\text{แรงสนาม}
\]

\begin{tcolorbox}[title=แรงสัมผัส]
แรงสัมผัส คือ แรงที่เกิดจากการสัมผัสกันโดยตรงระหว่างวัตถุ

ตัวอย่างเช่น แรงปกติ แรงตึงเชือก แรงเสียดทาน แรงผลัก และแรงดึง
\end{tcolorbox}

\begin{tcolorbox}[title=แรงสนาม]
แรงสนาม คือ แรงที่กระทำได้แม้วัตถุไม่สัมผัสกันโดยตรง

ตัวอย่างเช่น แรงโน้มถ่วง แรงไฟฟ้า และแรงแม่เหล็ก
\end{tcolorbox}

ในบทนี้จะเน้นแรงที่ใช้บ่อยในกลศาสตร์ ได้แก่

\[
\text{แรงโน้มถ่วง}
\]

\[
\text{แรงปกติ}
\]

\[
\text{แรงตึงเชือก}
\]

\[
\text{แรงเสียดทาน}
\]

\subsection{แรงลัพธ์}

ถ้ามีแรงหลายแรงกระทำต่อวัตถุ ผลรวมแบบเวกเตอร์ของแรงทั้งหมดเรียกว่า แรงลัพธ์

\[
\sum \vec{F}
\]

ถ้าแรงอยู่บนแกนเดียวกัน ให้ใช้เครื่องหมายบวกและลบตามทิศที่กำหนด

ถ้าแรงอยู่ในสองมิติ ให้แตกแรงเป็นองค์ประกอบตามแกนก่อน

\[
\sum F_x
\]

\[
\sum F_y
\]

\begin{tcolorbox}[title=ข้อควรจำ]
แรงลัพธ์ไม่ใช่ผลรวมของขนาดแรงอย่างเดียว แต่เป็นผลรวมแบบเวกเตอร์ ต้องคำนึงถึงทิศทางเสมอ
\end{tcolorbox}

\begin{examplebox}{ตัวอย่าง: แรงลัพธ์ในหนึ่งมิติ}
วัตถุถูกแรง $20\,\si{N}$ ไปทางขวา และแรง $8\,\si{N}$ ไปทางซ้าย จงหาแรงลัพธ์

เลือกขวาเป็นบวก

\[
\sum F=20-8
\]

\[
\sum F=12\,\si{N}
\]

ดังนั้นแรงลัพธ์มีขนาด

\[
12\,\si{N}
\]

และชี้ไปทางขวา
\end{examplebox}

\begin{exercisebox}{แบบฝึกหัด 5.1}
\begin{enumerate}
    \item จงจำแนกว่าแรงต่อไปนี้เป็นแรงสัมผัสหรือแรงสนาม: แรงเสียดทาน, แรงโน้มถ่วง, แรงตึงเชือก, แรงไฟฟ้า
    \item วัตถุถูกแรง $15\,\si{N}$ ไปทางขวา และ $6\,\si{N}$ ไปทางซ้าย จงหาแรงลัพธ์
    \item แรง $10\,\si{N}$ ไปทางเหนือ และแรง $10\,\si{N}$ ไปทางตะวันออก มีแรงลัพธ์ขนาดเท่าไร
    \item ถ้าแรงลัพธ์เป็นศูนย์ วัตถุต้องหยุดนิ่งเสมอหรือไม่
\end{enumerate}
\end{exercisebox}

% =========================================================
\section{กฎข้อที่หนึ่งของนิวตันและกรอบอ้างอิงเฉื่อย}

\subsection{กฎข้อที่หนึ่งของนิวตัน}

กฎข้อที่หนึ่งของนิวตันกล่าวถึงสภาพของวัตถุเมื่อไม่มีแรงลัพธ์กระทำ

\begin{tcolorbox}[title=กฎข้อที่หนึ่งของนิวตัน]
ถ้าแรงลัพธ์ที่กระทำต่อวัตถุเป็นศูนย์ วัตถุที่อยู่นิ่งจะยังคงอยู่นิ่ง และวัตถุที่เคลื่อนที่จะยังคงเคลื่อนที่ด้วยความเร็วคงตัวในแนวเส้นตรง
\[
\sum \vec{F}=0
\quad \Rightarrow \quad
\vec{a}=0
\]
\end{tcolorbox}

คำว่า

\[
\vec{a}=0
\]

หมายความว่า ความเร็วคงตัว ไม่ได้หมายความว่าวัตถุต้องหยุดนิ่งเสมอ

ดังนั้นสถานการณ์ที่แรงลัพธ์เป็นศูนย์มีได้สองแบบ

\[
\text{วัตถุอยู่นิ่ง}
\]

หรือ

\[
\text{วัตถุเคลื่อนที่ด้วยความเร็วคงตัวเป็นเส้นตรง}
\]

\subsection{ความเฉื่อย}

ความเฉื่อยคือแนวโน้มของวัตถุที่จะรักษาสภาพการเคลื่อนที่เดิม

ถ้าวัตถุอยู่นิ่ง มันมีแนวโน้มที่จะอยู่นิ่งต่อไป

ถ้าวัตถุกำลังเคลื่อนที่ด้วยความเร็วคงตัว มันมีแนวโน้มที่จะเคลื่อนที่แบบเดิมต่อไป

\begin{tcolorbox}[title=นิยาม]
ความเฉื่อย คือ สมบัติของวัตถุที่ต่อต้านการเปลี่ยนแปลงสภาพการเคลื่อนที่
\end{tcolorbox}

\subsection{กรอบอ้างอิงเฉื่อย}

กฎของนิวตันใช้ได้โดยตรงในกรอบอ้างอิงเฉื่อย

\begin{tcolorbox}[title=นิยาม]
กรอบอ้างอิงเฉื่อย คือ กรอบอ้างอิงที่วัตถุซึ่งไม่มีแรงลัพธ์กระทำจะไม่มีความเร่ง
\end{tcolorbox}

ตัวอย่างที่ดีโดยประมาณคือพื้นโลกในโจทย์ทั่วไปที่ไม่ต้องคำนึงถึงการหมุนของโลก

แต่รถที่กำลังเร่ง ลิฟต์ที่กำลังเร่ง หรือรถที่กำลังเลี้ยว ไม่ใช่กรอบอ้างอิงเฉื่อยอย่างแท้จริง

\begin{tcolorbox}[title=ข้อควรจำสำหรับ สอวน.]
เมื่อแก้โจทย์ด้วยกฎของนิวตัน ให้เลือกกรอบอ้างอิงที่เป็นกรอบเฉื่อยก่อน เช่น พื้นโลก ถ้าไปมองจากรถที่กำลังเร่ง ต้องระวังมากขึ้นเพราะจะมีแรงสมมติ ซึ่งยังไม่ใช่หัวข้อหลักของบทนี้
\end{tcolorbox}

\begin{examplebox}{ตัวอย่าง: สมดุลของวัตถุแขวน}
วัตถุมวล $m$ แขวนอยู่นิ่งด้วยเชือกเส้นเดียว จงหาแรงตึงเชือก

วัตถุอยู่นิ่ง ดังนั้น

\[
a=0
\]

จึงมี

\[
\sum F_y=0
\]

แรงที่กระทำต่อวัตถุมีสองแรงคือ แรงตึงเชือก $T$ ขึ้น และน้ำหนัก $mg$ ลง

เลือกขึ้นเป็นบวก

\[
T-mg=0
\]

ดังนั้น

\[
\boxed{T=mg}
\]
\end{examplebox}

\begin{exercisebox}{แบบฝึกหัด 5.2}
\begin{enumerate}
    \item วัตถุเคลื่อนที่เป็นเส้นตรงด้วยความเร็วคงตัว $5\,\si{m/s}$ แรงลัพธ์ที่กระทำต่อวัตถุเป็นเท่าไร
    \item กล่องวางนิ่งบนโต๊ะ มีแรงใดกระทำต่อกล่องบ้าง
    \item วัตถุถูกดึงขึ้นด้วยเชือกและอยู่นิ่ง จงเปรียบเทียบแรงตึงเชือกกับน้ำหนัก
    \item ถ้าลิฟต์กำลังเร่งขึ้น กรอบของลิฟต์เป็นกรอบเฉื่อยหรือไม่
\end{enumerate}
\end{exercisebox}

% =========================================================
\section{มวลและกฎข้อที่สองของนิวตัน}

\subsection{มวล}

มวลไม่ใช่น้ำหนัก มวลเป็นปริมาณที่บอกความเฉื่อยของวัตถุ

วัตถุที่มีมวลมาก เปลี่ยนสภาพการเคลื่อนที่ยากกว่า

วัตถุที่มีมวลน้อย เปลี่ยนสภาพการเคลื่อนที่ง่ายกว่า

\begin{tcolorbox}[title=นิยาม]
มวล คือ ปริมาณที่บอกความเฉื่อยของวัตถุ

หน่วยของมวลคือกิโลกรัม
\[
\si{kg}
\]
\end{tcolorbox}

\subsection{กฎข้อที่สองของนิวตัน}

กฎข้อที่สองของนิวตันเชื่อมแรงลัพธ์กับความเร่ง

\begin{tcolorbox}[title=กฎข้อที่สองของนิวตัน]
ความเร่งของวัตถุมีทิศเดียวกับแรงลัพธ์ และมีขนาดแปรผันตรงกับแรงลัพธ์ แต่แปรผกผันกับมวล

\[
\sum \vec{F}=m\vec{a}
\]
\end{tcolorbox}

ถ้าแยกตามแกน จะได้

\[
\sum F_x=ma_x
\]

\[
\sum F_y=ma_y
\]

\begin{tcolorbox}[title=ข้อควรจำ]
ใช้ $\sum \vec{F}=m\vec{a}$ กับแรงลัพธ์ ไม่ใช่แรงตัวใดตัวหนึ่ง

ถ้ามีแรงหลายแรง ต้องรวมแรงทั้งหมดก่อน
\end{tcolorbox}

\subsection{สมดุลเป็นกรณีพิเศษของกฎข้อที่สอง}

ถ้าวัตถุไม่มีความเร่ง

\[
\vec{a}=0
\]

จาก

\[
\sum \vec{F}=m\vec{a}
\]

จะได้

\[
\sum \vec{F}=0
\]

ดังนั้นสมดุลเป็นกรณีพิเศษของกฎข้อที่สอง

\[
\sum F_x=0
\]

\[
\sum F_y=0
\]

\begin{examplebox}{ตัวอย่าง: วัตถุถูกแรงลัพธ์คงตัว}
กล่องมวล $4\,\si{kg}$ อยู่บนพื้นลื่น ถูกแรงแนวนอน $20\,\si{N}$ จงหาความเร่ง

เพราะพื้นลื่น ไม่มีแรงเสียดทานในแนวราบ

\[
\sum F_x=ma_x
\]

\[
20=4a
\]

ดังนั้น

\[
\boxed{a=5\,\si{m/s^2}}
\]
\end{examplebox}

\begin{examplebox}{ตัวอย่าง: แรงสองแรงในแนวเดียวกัน}
วัตถุมวล $5\,\si{kg}$ ถูกแรง $30\,\si{N}$ ไปทางขวา และแรง $10\,\si{N}$ ไปทางซ้าย บนพื้นลื่น จงหาความเร่ง

เลือกขวาเป็นบวก

\[
\sum F_x=30-10
\]

\[
\sum F_x=20\,\si{N}
\]

ใช้กฎข้อที่สอง

\[
\sum F_x=ma_x
\]

\[
20=5a
\]

ดังนั้น

\[
\boxed{a=4\,\si{m/s^2}}
\]

ทิศไปทางขวา
\end{examplebox}

\begin{exercisebox}{แบบฝึกหัด 5.3}
\begin{enumerate}
    \item วัตถุมวล $2\,\si{kg}$ ถูกแรงลัพธ์ $10\,\si{N}$ จงหาความเร่ง
    \item วัตถุมวล $5\,\si{kg}$ มีความเร่ง $3\,\si{m/s^2}$ จงหาแรงลัพธ์
    \item วัตถุมวล $4\,\si{kg}$ ถูกแรง $25\,\si{N}$ ไปทางขวา และ $9\,\si{N}$ ไปทางซ้าย จงหาความเร่ง
    \item ถ้าแรงลัพธ์เพิ่มเป็นสองเท่า แต่มวลเท่าเดิม ความเร่งเปลี่ยนอย่างไร
    \item ถ้ามวลเพิ่มเป็นสองเท่า แต่แรงลัพธ์เท่าเดิม ความเร่งเปลี่ยนอย่างไร
\end{enumerate}
\end{exercisebox}

% =========================================================
\section{น้ำหนักและแรงโน้มถ่วง}

\subsection{แรงโน้มถ่วงใกล้ผิวโลก}

ใกล้ผิวโลก วัตถุมวล $m$ ถูกโลกดึงลงด้วยแรงโน้มถ่วง

\[
\vec{F}_g=m\vec{g}
\]

ถ้าเลือกแกนขึ้นเป็นบวก จะมี

\[
\vec{g}=-g\hat{j}
\]

ดังนั้น

\[
\vec{F}_g=-mg\hat{j}
\]

ขนาดของแรงโน้มถ่วงเรียกว่า น้ำหนัก

\[
F_g=mg
\]

\begin{tcolorbox}[title=นิยาม]
น้ำหนัก คือ ขนาดของแรงโน้มถ่วงที่โลกกระทำต่อวัตถุ

\[
W=mg
\]
\end{tcolorbox}

\subsection{มวลกับน้ำหนักต่างกันอย่างไร}

มวลเป็นสมบัติของวัตถุ และมีหน่วยเป็น

\[
\si{kg}
\]

น้ำหนักเป็นแรง และมีหน่วยเป็น

\[
\si{N}
\]

มวลของวัตถุเท่าเดิมไม่ว่าจะอยู่ที่ใด แต่น้ำหนักเปลี่ยนได้ถ้าค่า $g$ เปลี่ยน

\begin{center}
\begin{tabular}{|p{0.32\linewidth}|p{0.28\linewidth}|p{0.28\linewidth}|}
\hline
\textbf{หัวข้อ} & \textbf{มวล} & \textbf{น้ำหนัก} \\
\hline
ความหมาย & ความเฉื่อยของวัตถุ & แรงโน้มถ่วงที่กระทำต่อวัตถุ \\
\hline
ชนิดปริมาณ & สเกลาร์ & เวกเตอร์ \\
\hline
หน่วย & $\si{kg}$ & $\si{N}$ \\
\hline
สูตร & ไม่มีสูตรเดียวแบบแรง & $W=mg$ \\
\hline
เปลี่ยนตามสถานที่หรือไม่ & โดยทั่วไปไม่เปลี่ยน & เปลี่ยนตามค่า $g$ \\
\hline
\end{tabular}
\end{center}

\subsection{แรงปกติ}

แรงปกติคือแรงที่พื้นหรือผิวสัมผัสกระทำต่อวัตถุในทิศตั้งฉากกับผิวนั้น

\begin{tcolorbox}[title=นิยาม]
แรงปกติ $N$ คือแรงสัมผัสที่พื้นกระทำต่อวัตถุในทิศตั้งฉากกับผิวสัมผัส
\end{tcolorbox}

บนพื้นราบที่วัตถุไม่มีความเร่งในแนวดิ่ง จะได้

\[
N=mg
\]

แต่ไม่ควรจำว่า $N=mg$ เสมอ เพราะถ้ามีแรงอื่นในแนวดิ่ง หรือพื้นเอียง หรือระบบเร่ง ค่า $N$ อาจไม่เท่ากับ $mg$

\begin{examplebox}{ตัวอย่าง: กล่องบนพื้นราบ}
กล่องมวล $m$ วางนิ่งบนพื้นราบ จงหาแรงปกติ

แรงในแนวดิ่งมี

\[
N \text{ ขึ้น}
\]

และ

\[
mg \text{ ลง}
\]

เพราะกล่องไม่มีความเร่งในแนวดิ่ง

\[
\sum F_y=0
\]

เลือกขึ้นเป็นบวก

\[
N-mg=0
\]

ดังนั้น

\[
\boxed{N=mg}
\]
\end{examplebox}

\begin{examplebox}{ตัวอย่าง: แรงปกติในลิฟต์}
คนมวล $m$ ยืนบนเครื่องชั่งในลิฟต์ที่เร่งขึ้นด้วยความเร่ง $a$ จงหาแรงปกติ

แรงที่กระทำต่อคนคือ

\[
N \text{ ขึ้น}
\]

และ

\[
mg \text{ ลง}
\]

เลือกขึ้นเป็นบวก

\[
\sum F_y=ma
\]

\[
N-mg=ma
\]

ดังนั้น

\[
\boxed{N=m(g+a)}
\]

เครื่องชั่งจึงอ่านค่ามากกว่าน้ำหนักจริง
\end{examplebox}

\begin{exercisebox}{แบบฝึกหัด 5.4}
\begin{enumerate}
    \item วัตถุมวล $10\,\si{kg}$ อยู่บนโลก ใช้ $g=10\,\si{m/s^2}$ จงหาน้ำหนัก
    \item กล่องมวล $m$ วางบนพื้นราบนิ่ง จงหาแรงปกติ
    \item คนมวล $m$ อยู่ในลิฟต์ที่เร่งลงด้วยความเร่ง $a$ จงหาแรงปกติ
    \item ถ้าลิฟต์ตกอย่างอิสระ เครื่องชั่งในลิฟต์จะอ่านค่าเท่าไร
\end{enumerate}
\end{exercisebox}

% =========================================================
\section{กฎข้อที่สามของนิวตัน}

\subsection{คู่แรงกิริยาและปฏิกิริยา}

เมื่อวัตถุสองชิ้นมีปฏิสัมพันธ์กัน แรงจะเกิดเป็นคู่เสมอ

\begin{tcolorbox}[title=กฎข้อที่สามของนิวตัน]
ถ้าวัตถุ 1 ออกแรงกระทำต่อวัตถุ 2 วัตถุ 2 จะออกแรงกระทำต่อวัตถุ 1 ด้วยขนาดเท่ากัน แต่ทิศตรงข้าม

\[
\vec{F}_{12}=-\vec{F}_{21}
\]
\end{tcolorbox}

ตัวอย่างเช่น มือผลักกำแพง กำแพงก็ผลักมือกลับด้วยแรงขนาดเท่ากันในทิศตรงข้าม

\subsection{ข้อผิดพลาดที่พบบ่อย}

นักเรียนมักคิดว่าแรงกิริยาและปฏิกิริยาหักล้างกัน แต่จริง ๆ แล้วคู่แรงนี้กระทำบนคนละวัตถุ จึงไม่หักล้างกันในแผนภาพวัตถุอิสระของวัตถุเดียว

\begin{tcolorbox}[title=ข้อควรจำ]
คู่แรงกิริยาและปฏิกิริยามีลักษณะดังนี้

\[
\text{ขนาดเท่ากัน}
\]

\[
\text{ทิศตรงข้าม}
\]

\[
\text{เป็นแรงชนิดเดียวกัน}
\]

\[
\text{กระทำบนวัตถุคนละชิ้น}
\]
\end{tcolorbox}

\subsection{ตัวอย่างคู่แรง}

ถ้ากล่องวางบนโต๊ะ

แรงที่โลกกระทำต่อกล่องคือ

\[
\vec{F}_{\text{โลกต่อกล่อง}}
\]

คู่แรงของมันคือแรงที่กล่องกระทำต่อโลก

\[
\vec{F}_{\text{กล่องต่อโลก}}
\]

ถ้าโต๊ะดันกล่องขึ้นด้วยแรงปกติ

\[
\vec{N}_{\text{โต๊ะต่อกล่อง}}
\]

คู่แรงของมันคือแรงที่กล่องดันโต๊ะลง

\[
\vec{N}_{\text{กล่องต่อโต๊ะ}}
\]

\begin{exercisebox}{แบบฝึกหัด 5.5}
\begin{enumerate}
    \item คนผลักกำแพงไปทางขวา กำแพงออกแรงต่อคนในทิศใด
    \item กล่องวางบนโต๊ะ คู่แรงของแรงปกติที่โต๊ะกระทำต่อกล่องคือแรงใด
    \item คู่แรงกิริยาและปฏิกิริยาหักล้างกันหรือไม่ เพราะเหตุใด
    \item โลกดึงดวงจันทร์ และดวงจันทร์ดึงโลก แรงทั้งสองมีขนาดเท่ากันหรือไม่
\end{enumerate}
\end{exercisebox}

% =========================================================
\section{แผนภาพวัตถุอิสระ}

\subsection{ทำไมต้องวาดแผนภาพวัตถุอิสระ}

โจทย์กฎของนิวตันส่วนใหญ่ผิดเพราะนักเรียนลืมแรงบางแรง หรือใส่แรงที่ไม่ควรใส่ แผนภาพวัตถุอิสระช่วยให้เห็นชัดว่าแรงใดกระทำต่อวัตถุที่เราสนใจ

\begin{tcolorbox}[title=นิยาม]
แผนภาพวัตถุอิสระ หรือ free-body diagram คือรูปที่แสดงวัตถุหนึ่งชิ้นแยกออกจากสิ่งแวดล้อม และวาดเฉพาะแรงที่กระทำต่อวัตถุนั้น
\end{tcolorbox}

\subsection{ขั้นตอนการวาด}

\begin{enumerate}
    \item เลือกวัตถุที่ต้องการพิจารณา
    \item วาดวัตถุนั้นเป็นจุดหรือกล่องเล็ก ๆ
    \item วาดแรงภายนอกทุกแรงที่กระทำต่อวัตถุนั้น
    \item ไม่วาดแรงที่วัตถุนั้นกระทำต่อสิ่งอื่น
    \item เลือกแกนให้สะดวก
    \item แตกแรงเป็นองค์ประกอบตามแกน
    \item ตั้งสมการ $\sum F_x=ma_x$ และ $\sum F_y=ma_y$
\end{enumerate}

\begin{tcolorbox}[title=แรงที่พบบ่อยใน free-body diagram]
\[
mg \quad \text{น้ำหนัก ชี้ลง}
\]

\[
N \quad \text{แรงปกติ ตั้งฉากกับผิว}
\]

\[
T \quad \text{แรงตึงเชือก ตามแนวเชือก}
\]

\[
f \quad \text{แรงเสียดทาน ขนานกับผิว และต้านการไถลหรือแนวโน้มการไถล}
\]
\end{tcolorbox}

\subsection{หลักการเลือกแกน}

ถ้าวัตถุเคลื่อนที่บนพื้นราบ มักเลือกแกน $x$ ตามแนวพื้น และแกน $y$ ตั้งฉากกับพื้น

ถ้าวัตถุเคลื่อนที่บนพื้นเอียง มักเลือกแกนหนึ่งขนานพื้นเอียง และอีกแกนตั้งฉากกับพื้นเอียง

\begin{tcolorbox}[title=ข้อควรจำ]
การเลือกแกนที่ดีทำให้สมการง่ายขึ้น โดยเฉพาะพื้นเอียง ควรเลือกแกนขนานพื้นเอียงและตั้งฉากพื้นเอียง
\end{tcolorbox}

% =========================================================
\section{แรงตึงเชือกและระบบวัตถุหลายชิ้น}

\subsection{แรงตึงเชือก}

ในโจทย์พื้นฐาน ถ้าเชือกเบาและไม่ยืด และรอกเบาไร้แรงเสียดทาน แรงตึงเชือกจะมีขนาดเท่ากันตลอดเส้นเชือกเส้นเดียวกัน

\begin{tcolorbox}[title=แรงตึงเชือกในอุดมคติ]
ถ้าเชือกเบา ไม่ยืด และรอกไร้แรงเสียดทาน

\[
T=\text{ค่าคงตัวตลอดเชือกเส้นเดียวกัน}
\]
\end{tcolorbox}

\subsection{ระบบสองมวลบนพื้นลื่น}

มวล $m_1$ และ $m_2$ วางบนพื้นลื่น เชื่อมด้วยเชือกเบา และถูกดึงด้วยแรงภายนอก $F$ ที่มวล $m_2$

ถ้าพิจารณาทั้งระบบเป็นวัตถุเดียว มวลรวมคือ

\[
m_1+m_2
\]

แรงภายนอกในแนวราบคือ

\[
F
\]

ดังนั้น

\[
F=(m_1+m_2)a
\]

\[
a=\frac{F}{m_1+m_2}
\]

ถ้าต้องการหาแรงตึงเชือก ให้พิจารณาเฉพาะมวล $m_1$

แรงในแนวราบที่กระทำต่อ $m_1$ คือ $T$

\[
T=m_1a
\]

ดังนั้น

\[
T=m_1\left(\frac{F}{m_1+m_2}\right)
\]

\[
T=\frac{m_1F}{m_1+m_2}
\]

\begin{examplebox}{ตัวอย่าง: ระบบสองมวลบนพื้นลื่น}
มวล $2\,\si{kg}$ และ $3\,\si{kg}$ เชื่อมด้วยเชือกเบา วางบนพื้นลื่น และถูกดึงด้วยแรง $20\,\si{N}$ จงหาความเร่งของระบบ

พิจารณาทั้งระบบ

\[
F=(m_1+m_2)a
\]

แทนค่า

\[
20=(2+3)a
\]

\[
20=5a
\]

ดังนั้น

\[
\boxed{a=4\,\si{m/s^2}}
\]
\end{examplebox}

\begin{exercisebox}{แบบฝึกหัด 5.6}
\begin{enumerate}
    \item มวล $2\,\si{kg}$ และ $6\,\si{kg}$ เชื่อมด้วยเชือกเบา วางบนพื้นลื่น ถูกดึงด้วยแรง $16\,\si{N}$ จงหาความเร่ง
    \item จากข้อก่อนหน้า จงหาแรงตึงเชือก
    \item มวล $m_1$ และ $m_2$ วางบนพื้นลื่นและถูกดึงด้วยแรง $F$ จงพิสูจน์ว่า $a=F/(m_1+m_2)$
\end{enumerate}
\end{exercisebox}

% =========================================================
\section{พื้นเอียง}

\subsection{ทำไมต้องแตกน้ำหนัก}

เมื่อวัตถุอยู่บนพื้นเอียง น้ำหนัก $mg$ ชี้ลงแนวดิ่ง แต่การเคลื่อนที่มักเกิดตามแนวพื้นเอียง เราจึงแตกน้ำหนักเป็นสององค์ประกอบ

\[
mg\sin\theta
\]

ขนานพื้นเอียง และ

\[
mg\cos\theta
\]

ตั้งฉากกับพื้นเอียง

โดย $\theta$ คือมุมของพื้นเอียงกับแนวระดับ

\begin{tcolorbox}[title=องค์ประกอบของน้ำหนักบนพื้นเอียง]
สำหรับพื้นเอียงมุม $\theta$

องค์ประกอบขนานพื้นเอียง

\[
mg\sin\theta
\]

องค์ประกอบตั้งฉากกับพื้นเอียง

\[
mg\cos\theta
\]
\end{tcolorbox}

\subsection{พื้นเอียงลื่น}

ถ้าพื้นเอียงลื่น ไม่มีแรงเสียดทาน

แรงตามแนวพื้นเอียงคือ

\[
mg\sin\theta
\]

แรงตั้งฉากพื้นเอียงสมดุลกัน

\[
N=mg\cos\theta
\]

ตามแนวพื้นเอียง

\[
\sum F=ma
\]

\[
mg\sin\theta=ma
\]

ดังนั้น

\[
a=g\sin\theta
\]

\begin{tcolorbox}[title=พื้นเอียงลื่น]
\[
a=g\sin\theta
\]

\[
N=mg\cos\theta
\]
\end{tcolorbox}

\begin{examplebox}{ตัวอย่าง: วัตถุไถลบนพื้นเอียงลื่น}
วัตถุไถลลงพื้นเอียงลื่นมุม $30^\circ$ จงหาความเร่ง

ใช้

\[
a=g\sin\theta
\]

แทนค่า

\[
a=g\sin30^\circ
\]

\[
a=\frac{g}{2}
\]

ถ้าใช้ $g=10\,\si{m/s^2}$

\[
\boxed{a=5\,\si{m/s^2}}
\]
\end{examplebox}

\begin{exercisebox}{แบบฝึกหัด 5.7}
\begin{enumerate}
    \item วัตถุมวล $m$ อยู่บนพื้นเอียงลื่นมุม $30^\circ$ จงหา $a$
    \item วัตถุมวล $m$ อยู่บนพื้นเอียงลื่นมุม $\theta$ จงหา $N$
    \item พื้นเอียงมุมใดทำให้วัตถุไถลลงด้วยความเร่ง $g/2$
    \item ถ้ามุมพื้นเอียงเพิ่มขึ้น ความเร่งของวัตถุบนพื้นลื่นเปลี่ยนอย่างไร
\end{enumerate}
\end{exercisebox}

% =========================================================
\section{แรงเสียดทาน}

\subsection{แนวคิดของแรงเสียดทาน}

แรงเสียดทานเป็นแรงสัมผัสที่เกิดขึ้นระหว่างผิวสองผิว และมีทิศขนานกับผิวสัมผัส

แรงเสียดทานมีหน้าที่ต้านการไถล หรือแนวโน้มที่จะไถล

\begin{tcolorbox}[title=นิยาม]
แรงเสียดทาน คือ แรงที่ผิวสัมผัสกระทำต่อวัตถุในทิศขนานกับผิว และต้านการไถลหรือแนวโน้มการไถล
\end{tcolorbox}

แรงเสียดทานแบ่งเป็นสองชนิดหลัก

\[
\text{แรงเสียดทานสถิต}
\]

และ

\[
\text{แรงเสียดทานจลน์}
\]

\subsection{แรงเสียดทานสถิต}

แรงเสียดทานสถิตเกิดเมื่อผิวสัมผัสยังไม่ไถลผ่านกัน

แรงเสียดทานสถิตจะปรับค่าตามความจำเป็นจนถึงค่าสูงสุด

\[
f_s\leq \mu_s N
\]

ค่าสูงสุดคือ

\[
f_{s,\max}=\mu_s N
\]

\begin{tcolorbox}[title=แรงเสียดทานสถิต]
\[
f_s\leq \mu_s N
\]

\[
f_{s,\max}=\mu_s N
\]
\end{tcolorbox}

\begin{tcolorbox}[title=ข้อควรระวัง]
ห้ามเขียน $f_s=\mu_s N$ ทันทีทุกครั้ง เพราะแรงเสียดทานสถิตไม่ได้มีค่าเท่าค่าสูงสุดเสมอไป

เขียน $f_s=\mu_s N$ ได้เมื่อวัตถุกำลังจะเริ่มไถล
\end{tcolorbox}

\subsection{แรงเสียดทานจลน์}

แรงเสียดทานจลน์เกิดเมื่อผิวสัมผัสไถลผ่านกันแล้ว

ขนาดของแรงเสียดทานจลน์คือ

\[
f_k=\mu_k N
\]

\begin{tcolorbox}[title=แรงเสียดทานจลน์]
\[
f_k=\mu_k N
\]
\end{tcolorbox}

โดยทั่วไป

\[
\mu_s>\mu_k
\]

หมายความว่า การทำให้วัตถุเริ่มไถลมักยากกว่าการทำให้วัตถุไถลต่อไป

\subsection{พื้นเอียงที่มีแรงเสียดทาน}

วัตถุมวล $m$ ไถลลงพื้นเอียงมุม $\theta$ โดยมีแรงเสียดทานจลน์

แรงตามแนวพื้นเอียงลงคือ

\[
mg\sin\theta
\]

แรงเสียดทานจลน์ต้านการไถลขึ้นตามพื้นเอียง

\[
f_k=\mu_k N
\]

และ

\[
N=mg\cos\theta
\]

ดังนั้น

\[
f_k=\mu_k mg\cos\theta
\]

ตั้งสมการตามแนวพื้นเอียงลงเป็นบวก

\[
mg\sin\theta-\mu_k mg\cos\theta=ma
\]

ตัด $m$

\[
a=g(\sin\theta-\mu_k\cos\theta)
\]

\begin{tcolorbox}[title=พื้นเอียงมีแรงเสียดทานจลน์]
ถ้าวัตถุไถลลงพื้นเอียง

\[
a=g(\sin\theta-\mu_k\cos\theta)
\]
\end{tcolorbox}

\begin{examplebox}{ตัวอย่าง: พื้นราบมีแรงเสียดทาน}
กล่องมวล $10\,\si{kg}$ ถูกดึงด้วยแรง $60\,\si{N}$ บนพื้นราบ โดยมี $\mu_k=0.2$ ใช้ $g=10\,\si{m/s^2}$ จงหาความเร่ง

บนพื้นราบ

\[
N=mg=10(10)=100\,\si{N}
\]

แรงเสียดทานจลน์คือ

\[
f_k=\mu_k N
\]

\[
f_k=0.2(100)=20\,\si{N}
\]

แรงลัพธ์แนวราบคือ

\[
\sum F_x=60-20=40\,\si{N}
\]

ใช้

\[
\sum F_x=ma
\]

\[
40=10a
\]

ดังนั้น

\[
\boxed{a=4\,\si{m/s^2}}
\]
\end{examplebox}

\begin{examplebox}{ตัวอย่าง: มุมเริ่มไถลบนพื้นเอียง}
วัตถุวางบนพื้นเอียงที่ค่อย ๆ เพิ่มมุม ถ้าสัมประสิทธิ์เสียดทานสถิตเป็น $\mu_s$ จงหามุมที่วัตถุเริ่มไถล

ที่กำลังจะเริ่มไถล

\[
f_s=f_{s,\max}=\mu_s N
\]

แรงตามแนวพื้นเอียงลงคือ

\[
mg\sin\theta
\]

แรงเสียดทานสถิตต้านขึ้นตามพื้นเอียง

\[
f_s=\mu_s N
\]

แรงตั้งฉากพื้นเอียงให้

\[
N=mg\cos\theta
\]

ที่จุดเริ่มไถล

\[
mg\sin\theta=\mu_s mg\cos\theta
\]

ตัด $mg$

\[
\sin\theta=\mu_s\cos\theta
\]

ดังนั้น

\[
\tan\theta=\mu_s
\]

\[
\boxed{\theta=\tan^{-1}\mu_s}
\]
\end{examplebox}

\begin{exercisebox}{แบบฝึกหัด 5.8}
\begin{enumerate}
    \item กล่องมวล $5\,\si{kg}$ วางบนพื้นราบ มี $\mu_k=0.3$ ใช้ $g=10\,\si{m/s^2}$ จงหาแรงเสียดทานจลน์
    \item กล่องมวล $10\,\si{kg}$ ถูกดึงด้วยแรง $50\,\si{N}$ บนพื้นราบ มี $\mu_k=0.2$ ใช้ $g=10\,\si{m/s^2}$ จงหาความเร่ง
    \item วัตถุอยู่บนพื้นเอียงมุม $30^\circ$ และมี $\mu_k=0.1$ จงเขียนสมการความเร่งลงตามพื้นเอียง
    \item วัตถุเริ่มไถลบนพื้นเอียงเมื่อมุมเท่ากับ $37^\circ$ จงหา $\mu_s$ โดยใช้ $\sin37^\circ=3/5$ และ $\cos37^\circ=4/5$
    \item อธิบายความแตกต่างระหว่าง $f_s\leq \mu_sN$ และ $f_k=\mu_kN$
\end{enumerate}
\end{exercisebox}

% =========================================================
\section{แบบจำลองการวิเคราะห์ด้วยกฎข้อที่สองของนิวตัน}


\begin{tcolorbox}[title=แบบจำลองที่ 1: อนุภาคสมดุล]
ถ้าวัตถุอยู่นิ่ง หรือเคลื่อนที่ด้วยความเร็วคงตัวเป็นเส้นตรง

\[
\vec{a}=0
\]

ดังนั้น

\[
\sum \vec{F}=0
\]

แยกแกนได้เป็น

\[
\sum F_x=0
\]

\[
\sum F_y=0
\]
\end{tcolorbox}

\begin{tcolorbox}[title=แบบจำลองที่ 2: อนุภาคภายใต้แรงลัพธ์]
ถ้าวัตถุมีความเร่ง

\[
\vec{a}\neq 0
\]

ต้องใช้

\[
\sum \vec{F}=m\vec{a}
\]

แยกแกนได้เป็น

\[
\sum F_x=ma_x
\]

\[
\sum F_y=ma_y
\]
\end{tcolorbox}

\subsection{แรงตึงเชือกและรอกอุดมคติ}

ในโจทย์พื้นฐาน ถ้าเชือกเบา ไม่ยืด และรอกเบาไร้แรงเสียดทาน จะถือว่าแรงตึงเชือกมีขนาดเท่ากันตลอดเชือกเส้นเดียวกัน

\[
T=\text{ค่าคงตัวตลอดเชือกเส้นเดียวกัน}
\]

แต่ทิศของแรงตึงจะเปลี่ยนไปตามแนวเชือก

\begin{tcolorbox}[title=ข้อควรจำ]
แรงตึงเชือกเป็นแรงที่เชือกดึงวัตถุ ดังนั้นแรงตึงเชือกจะชี้ออกจากวัตถุไปตามแนวเชือกเสมอ
\end{tcolorbox}

\subsection{เครื่องชั่งสปริงอ่านค่าอะไร}

เครื่องชั่งสปริงที่ต่ออยู่ในเชือกอ่านค่าแรงตึงเชือก ไม่ได้อ่านค่าแรงลัพธ์ของระบบ

ถ้าเชือกเบาและเครื่องชั่งเบา ค่าอ่านของเครื่องชั่งคือ

\[
T
\]

\begin{tcolorbox}[title=ข้อควรระวัง]
ในระบบหลายมวล อย่าหาแรงลัพธ์ทั้งระบบแล้วตอบเป็นค่าเครื่องชั่งทันที ต้องพิจารณาว่าเครื่องชั่งอยู่ตรงส่วนใดของเชือก และแรงตึงเชือกส่วนนั้นมีค่าเท่าไร
\end{tcolorbox}

\subsection{ระบบรอกแบบ Atwood}

ให้มวล $m_1$ และ $m_2$ ต่อด้วยเชือกเบาผ่านรอกเบาไร้แรงเสียดทาน และให้

\[
m_1>m_2
\]

ดังนั้น $m_1$ ลง และ $m_2$ ขึ้น ด้วยความเร่งขนาดเดียวกัน $a$

สำหรับมวล $m_1$ เลือกลงเป็นบวก

\[
m_1g-T=m_1a
\]

สำหรับมวล $m_2$ เลือกขึ้นเป็นบวก

\[
T-m_2g=m_2a
\]

บวกสองสมการ

\[
m_1g-m_2g=(m_1+m_2)a
\]

ดังนั้น

\[
a=\frac{(m_1-m_2)g}{m_1+m_2}
\]

\begin{tcolorbox}[title=Atwood machine]
ถ้า $m_1>m_2$

\[
a=\frac{(m_1-m_2)g}{m_1+m_2}
\]

และหาแรงตึงเชือกได้จาก

\[
T=m_2(g+a)
\]

หรือ

\[
T=m_1(g-a)
\]
\end{tcolorbox}

\begin{examplebox}{ตัวอย่าง: ระบบรอกพื้นฐาน}
มวล $6\,\si{kg}$ และ $4\,\si{kg}$ ต่อด้วยเชือกเบาผ่านรอกเบาไร้แรงเสียดทาน จงหาความเร่งของระบบ ใช้ $g=10\,\si{m/s^2}$

เพราะมวล $6\,\si{kg}$ มากกว่า จึงเคลื่อนที่ลง

\[
a=\frac{(m_1-m_2)g}{m_1+m_2}
\]

แทนค่า

\[
a=\frac{(6-4)(10)}{6+4}
\]

\[
a=2\,\si{m/s^2}
\]
\end{examplebox}

\begin{exercisebox}{แบบฝึกหัดเพิ่มเติม: ระบบเชือกและรอก}
\begin{enumerate}
    \item มวล $5\,\si{kg}$ และ $3\,\si{kg}$ ต่อด้วยเชือกเบาผ่านรอกเบาไร้แรงเสียดทาน จงหาความเร่งของระบบ
    \item จากข้อก่อนหน้า จงหาแรงตึงเชือก ใช้ $g=10\,\si{m/s^2}$
    \item มวล $m_1$ วางบนโต๊ะลื่น ต่อเชือกผ่านรอกกับมวลแขวน $m_2$ จงพิสูจน์ว่า
    \[
    a=\frac{m_2g}{m_1+m_2}
    \]
    \item เครื่องชั่งสปริงอยู่ระหว่างมวลสองก้อนบนพื้นลื่น จงอธิบายว่าค่าอ่านของเครื่องชั่งเท่ากับแรงใด
\end{enumerate}
\end{exercisebox}

% =========================================================
\section{วิธีแก้โจทย์กฎของนิวตัน}

\subsection{ขั้นตอนมาตรฐาน}

\begin{tcolorbox}[title=ขั้นตอนแก้โจทย์กฎของนิวตัน]
\begin{enumerate}
    \item อ่านโจทย์และเลือกวัตถุที่จะพิจารณา
    \item วาดแผนภาพวัตถุอิสระ
    \item เลือกแกนให้สะดวก
    \item แตกแรงที่เอียงออกเป็นองค์ประกอบ
    \item เขียนสมการ
    \[
    \sum F_x=ma_x
    \]
    และ
    \[
    \sum F_y=ma_y
    \]
    \item ใส่เงื่อนไขพิเศษ เช่น วัตถุอยู่นิ่ง วัตถุเคลื่อนที่ด้วยความเร็วคงตัว หรือวัตถุหลายชิ้นมีความเร่งเท่ากัน
    \item แก้สมการและตรวจทิศทางของคำตอบ
\end{enumerate}
\end{tcolorbox}

\subsection{สิ่งที่ต้องระวัง}

\begin{enumerate}
    \item อย่าใส่แรง ``$ma$'' ลงใน free-body diagram เพราะ $ma$ ไม่ใช่แรง แต่เป็นผลของแรงลัพธ์
    \item อย่าคิดว่าแรงปกติเท่ากับ $mg$ เสมอ
    \item อย่าใช้ $f_s=\mu_sN$ ถ้าวัตถุยังไม่ได้อยู่ที่จุดใกล้ไถล
    \item คู่แรงกิริยาและปฏิกิริยาไม่อยู่ใน free-body diagram ของวัตถุเดียวกัน
    \item ต้องแตกแรงตามแกนก่อนใช้สมการ
\end{enumerate}

\begin{examplebox}{ตัวอย่างรวม: แรงดึงทำมุมกับพื้นราบ}
กล่องมวล $m$ บนพื้นลื่นถูกดึงด้วยแรง $F$ ทำมุม $\theta$ เหนือแนวระดับ จงหาความเร่งในแนวราบและแรงปกติ

แตกแรงดึง

\[
F_x=F\cos\theta
\]

\[
F_y=F\sin\theta
\]

แกน $x$

\[
\sum F_x=ma_x
\]

\[
F\cos\theta=ma
\]

ดังนั้น

\[
\boxed{a=\frac{F\cos\theta}{m}}
\]

แกน $y$ วัตถุไม่มีความเร่งในแนวดิ่ง

\[
\sum F_y=0
\]

แรงในแนวดิ่งมี $N$ ขึ้น, $F\sin\theta$ ขึ้น, และ $mg$ ลง

\[
N+F\sin\theta-mg=0
\]

ดังนั้น

\[
\boxed{N=mg-F\sin\theta}
\]
\end{examplebox}

% =========================================================
\section{ชุดโจทย์รวมท้ายบท: ระดับผสมสำหรับ สอวน. ค่าย 1}

\subsection*{ระดับพื้นฐาน}

\begin{enumerate}
    \item วัตถุมวล $3\,\si{kg}$ ถูกแรงลัพธ์ $18\,\si{N}$ จงหาความเร่ง
    \item วัตถุมวล $5\,\si{kg}$ วางบนพื้นราบ ใช้ $g=10\,\si{m/s^2}$ จงหาแรงปกติ
    \item วัตถุมวล $m$ แขวนอยู่นิ่งด้วยเชือกเส้นเดียว จงหาแรงตึงเชือก
    \item กล่องบนพื้นลื่นถูกแรง $F$ ดึงในแนวราบ จงหาความเร่ง
\end{enumerate}

\subsection*{ระดับกลาง}

\begin{enumerate}
    \setcounter{enumi}{4}
    \item วัตถุมวล $m$ ไถลลงพื้นเอียงลื่นมุม $\theta$ จงพิสูจน์ว่า $a=g\sin\theta$
    \item กล่องมวล $m$ บนพื้นราบถูกดึงด้วยแรง $F$ ทำมุม $\theta$ เหนือแนวระดับ จงหา $N$ และ $a$
    \item วัตถุวางบนพื้นเอียงมุม $\theta$ และกำลังจะเริ่มไถล จงพิสูจน์ว่า $\mu_s=\tan\theta$
    \item มวล $m_1$ และ $m_2$ วางบนพื้นลื่น เชื่อมด้วยเชือก และถูกดึงด้วยแรง $F$ ที่มวล $m_2$ จงหา $a$ และ $T$
\end{enumerate}

\subsection*{ระดับยาก}

\begin{enumerate}
    \setcounter{enumi}{8}
    \item วัตถุมวล $m$ ไถลลงพื้นเอียงมุม $\theta$ โดยมีแรงเสียดทานจลน์ $\mu_k$ จงหาความเร่ง
    \item คนมวล $m$ ยืนบนเครื่องชั่งในลิฟต์ที่เร่งขึ้นด้วยความเร่ง $a$ จงหาแรงปกติ
    \item กล่องมวล $m$ บนพื้นราบมี $\mu_s$ ถูกดึงด้วยแรง $F$ ทำมุม $\theta$ เหนือแนวระดับ จงหาเงื่อนไขที่กล่องเริ่มเคลื่อนที่
    \item มวล $m$ ถูกดันติดกับผนังแนวดิ่งด้วยแรงแนวนอน $F$ ถ้าสัมประสิทธิ์เสียดทานสถิตระหว่างมวลกับผนังคือ $\mu_s$ จงหาแรง $F$ น้อยที่สุดที่ทำให้มวลไม่ไถลลง
\end{enumerate}

% =========================================================
\section{เฉลยย่อท้ายบท}

\begin{enumerate}
    \item
    \[
    \sum F=ma
    \]
    \[
    18=3a
    \]
    \[
    a=6\,\si{m/s^2}
    \]

    \item
    \[
    N=mg=5(10)=50\,\si{N}
    \]

    \item
    อยู่นิ่งจึงมี
    \[
    T-mg=0
    \]
    \[
    T=mg
    \]

    \item
    \[
    F=ma
    \]
    \[
    a=\frac{F}{m}
    \]

    \item
    เลือกแกนขนานพื้นเอียง
    \[
    \sum F=ma
    \]
    \[
    mg\sin\theta=ma
    \]
    \[
    a=g\sin\theta
    \]

    \item
    แตกแรง
    \[
    F_x=F\cos\theta
    \]
    \[
    F_y=F\sin\theta
    \]
    แนวราบ
    \[
    F\cos\theta=ma
    \]
    \[
    a=\frac{F\cos\theta}{m}
    \]
    แนวดิ่ง
    \[
    N+F\sin\theta-mg=0
    \]
    \[
    N=mg-F\sin\theta
    \]

    \item
    ที่กำลังจะเริ่มไถล
    \[
    f_s=\mu_sN
    \]
    และ
    \[
    mg\sin\theta=\mu_smg\cos\theta
    \]
    ดังนั้น
    \[
    \mu_s=\tan\theta
    \]

    \item
    พิจารณาทั้งระบบ
    \[
    F=(m_1+m_2)a
    \]
    \[
    a=\frac{F}{m_1+m_2}
    \]
    พิจารณา $m_1$
    \[
    T=m_1a
    \]
    \[
    T=\frac{m_1F}{m_1+m_2}
    \]

    \item
    แรงลงตามพื้นเอียงคือ
    \[
    mg\sin\theta
    \]
    แรงเสียดทานจลน์ขึ้นตามพื้นเอียงคือ
    \[
    f_k=\mu_kmg\cos\theta
    \]
    ดังนั้น
    \[
    mg\sin\theta-\mu_kmg\cos\theta=ma
    \]
    \[
    a=g(\sin\theta-\mu_k\cos\theta)
    \]

    \item
    เลือกขึ้นเป็นบวก
    \[
    N-mg=ma
    \]
    \[
    N=m(g+a)
    \]

    \item
    ที่เริ่มเคลื่อนที่
    \[
    F\cos\theta=f_{s,\max}
    \]
    โดย
    \[
    f_{s,\max}=\mu_sN
    \]
    และ
    \[
    N=mg-F\sin\theta
    \]
    ดังนั้น
    \[
    F\cos\theta=\mu_s(mg-F\sin\theta)
    \]
    \[
    F(\cos\theta+\mu_s\sin\theta)=\mu_smg
    \]
    \[
    F=\frac{\mu_smg}{\cos\theta+\mu_s\sin\theta}
    \]

    \item
    แรงปกติจากผนังคือ
    \[
    N=F
    \]
    แรงเสียดทานสถิตสูงสุดคือ
    \[
    f_{s,\max}=\mu_sF
    \]
    เพื่อไม่ให้มวลไถลลง ต้องมี
    \[
    \mu_sF\geq mg
    \]
    ดังนั้น
    \[
    F_{\min}=\frac{mg}{\mu_s}
    \]
\end{enumerate}

% =========================================================
\section{สรุปสูตรสำคัญ}

\begin{tcolorbox}[title=สรุปบทที่ 5: กฎการเคลื่อนที่]
แรงลัพธ์

\[
\sum \vec{F}
\]

กฎข้อที่หนึ่งของนิวตัน

\[
\sum \vec{F}=0
\quad \Rightarrow \quad
\vec{a}=0
\]

กฎข้อที่สองของนิวตัน

\[
\sum \vec{F}=m\vec{a}
\]

แยกแกน

\[
\sum F_x=ma_x
\]

\[
\sum F_y=ma_y
\]

น้ำหนัก

\[
W=mg
\]

แรงปกติบนพื้นราบเมื่อไม่มีความเร่งแนวดิ่ง

\[
N=mg
\]

พื้นเอียงลื่น

\[
N=mg\cos\theta
\]

\[
a=g\sin\theta
\]
\end{tcolorbox}

\begin{tcolorbox}[title=สรุปบทที่ 5: กฎการเคลื่อนที่ ต่อ]
กฎข้อที่สามของนิวตัน

\[
\vec{F}_{12}=-\vec{F}_{21}
\]

แรงเสียดทานสถิต

\[
f_s\leq \mu_sN
\]

\[
f_{s,\max}=\mu_sN
\]

แรงเสียดทานจลน์

\[
f_k=\mu_kN
\]

พื้นเอียงมีแรงเสียดทานจลน์และวัตถุไถลลง

\[
a=g(\sin\theta-\mu_k\cos\theta)
\]

มุมเริ่มไถล

\[
\tan\theta=\mu_s
\]
\end{tcolorbox}

% =========================================================
% =========================================================
\section{ข้อสอบจริง สอวน. ที่ใช้ความรู้บทกฎการเคลื่อนที่}

% =========================================================
\subsection{ปี 2560 ข้อ 2: กระดานลื่นที่ความชันค่อย ๆ ลดลง}

\vspace{1.5em}

\IfFileExists{img/posn60q2.png}{
\begin{center}
    \includegraphics[width=1\linewidth]{img/posn60q2.png}
\end{center}
}{
\begin{tcolorbox}[title=ตำแหน่งภาพที่แนะนำ]
ให้ตัดภาพข้อสอบปี 2560 ข้อ 2 แล้วบันทึกเป็น

\[
\texttt{img/posn60q2.png}
\]
\end{tcolorbox}
}

\vspace{1em}

\begin{examplebox}{เฉลยแนวคิด}
กระดานลื่นช่วงบนชันมาก แล้วค่อย ๆ ลดความชันลงที่ปลายล่าง

แรงที่ทำให้เด็กเร่งลงตามกระดานคือองค์ประกอบของน้ำหนักตามแนวกระดาน

\[
mg\sin\theta
\]

ดังนั้นความเร่งตามแนวกระดานคือ

\[
a=g\sin\theta
\]

เมื่อกระดานค่อย ๆ ลดความชันลง มุม $\theta$ ลดลง ทำให้

\[
\sin\theta \text{ ลดลง}
\]

ดังนั้น

\[
a \text{ ลดลง}
\]

แต่ขณะที่เด็กยังเคลื่อนที่ลงและยังมีความเร่งไปตามทิศการเคลื่อนที่ ความเร็วของเด็กยังเพิ่มขึ้น

สรุปคือ

\[
\boxed{\text{ขนาดของความเร็วเพิ่มขึ้น แต่ขนาดของความเร่งลดลง}}
\]

\[
\boxed{\text{ตอบ D}}
\]
\end{examplebox}

% =========================================================
\subsection{ปี 2561 ข้อ 2: ลิงกับมวลถ่วงผ่านรอก}

\vspace{1.5em}

\IfFileExists{img/posn61q2.png}{
\begin{center}
    \includegraphics[width=1\linewidth]{img/posn61q2.png}
\end{center}
}{
\begin{tcolorbox}[title=ตำแหน่งภาพที่แนะนำ]
ให้ตัดภาพข้อสอบปี 2561 ข้อ 2 แล้วบันทึกเป็น

\[
\texttt{img/posn61q2.png}
\]
\end{tcolorbox}
}

\vspace{1em}

\begin{examplebox}{เฉลยละเอียด}
ระบบเป็นมวลสองก้อนต่อด้วยเชือกเบาผ่านรอกเบาไร้แรงเสียดทาน

มวลของลิงคือ

\[
m_1=64\,\si{kg}
\]

มวลอีกด้านคือ

\[
m_2=34\,\si{kg}
\]

เพราะ

\[
m_1>m_2
\]

ลิงจึงเคลื่อนที่ลง และมวลอีกด้านเคลื่อนที่ขึ้น

ความเร่งของระบบ Atwood คือ

\[
a=\frac{(m_1-m_2)g}{m_1+m_2}
\]

แทนค่า

\[
a=\frac{(64-34)g}{64+34}
\]

\[
a=\frac{30g}{98}
\]

ใช้

\[
g=9.8\,\si{m/s^2}
\]

จะได้

\[
a=3.0\,\si{m/s^2}
\]

ลิงเริ่มจากหยุดนิ่งและเคลื่อนที่ลงเป็นระยะ

\[
s=6.0\,\si{m}
\]

ใช้สมการ

\[
v^2=u^2+2as
\]

โดย

\[
u=0
\]

ดังนั้น

\[
v^2=2(3.0)(6.0)
\]

\[
v^2=36
\]

\[
v=6.0\,\si{m/s}
\]

\[
\boxed{v=6.0\,\si{m/s}}
\]
\end{examplebox}

% =========================================================
\subsection{ปี 2561 ข้อ 4: เครื่องชั่งสปริงระหว่างมวลสองก้อน}

\vspace{1.5em}

\IfFileExists{img/posn61q4.png}{
\begin{center}
    \includegraphics[width=1\linewidth]{img/posn61q4.png}
\end{center}
}{
\begin{tcolorbox}[title=ตำแหน่งภาพที่แนะนำ]
ให้ตัดภาพข้อสอบปี 2561 ข้อ 4 แล้วบันทึกเป็น

\[
\texttt{img/posn61q4.png}
\]
\end{tcolorbox}
}

\vspace{1em}

\begin{examplebox}{เฉลยละเอียด}
มีมวลสองก้อนบนพื้นลื่น

\[
m_1=3.0\,\si{kg}
\]

\[
m_2=2.0\,\si{kg}
\]

มีแรงภายนอก $5.0\,\si{N}$ กระทำต่อมวลซ้ายไปทางขวา และแรง $15\,\si{N}$ กระทำต่อมวลขวาไปทางขวา

พิจารณาทั้งระบบ

\[
\sum F=(m_1+m_2)a
\]

\[
5.0+15=(3.0+2.0)a
\]

\[
20=5a
\]

ดังนั้น

\[
a=4.0\,\si{m/s^2}
\]

เครื่องชั่งสปริงอ่านค่าแรงตึง $T$ ระหว่างมวลสองก้อน

พิจารณามวล $3.0\,\si{kg}$

แรงในแนวราบที่กระทำต่อมวลนี้คือ

\[
5.0\,\si{N}
\]

และแรงตึงจากเครื่องชั่ง

\[
T
\]

ทั้งสองแรงชี้ไปทางขวา

ใช้

\[
\sum F=ma
\]

\[
5.0+T=(3.0)(4.0)
\]

\[
5.0+T=12
\]

ดังนั้น

\[
T=7.0\,\si{N}
\]

\[
\boxed{7.0\,\si{N}}
\]
\end{examplebox}

% =========================================================
\subsection{ปี 2561 ข้อ 5: มวลบนโต๊ะลื่นต่อกับมวลแขวน}

\vspace{1.5em}

\IfFileExists{img/posn61q5.png}{
\begin{center}
    \includegraphics[width=1\linewidth]{img/posn61q5.png}
\end{center}
}{
\begin{tcolorbox}[title=ตำแหน่งภาพที่แนะนำ]
ให้ตัดภาพข้อสอบปี 2561 ข้อ 5 แล้วบันทึกเป็น

\[
\texttt{img/posn61q5.png}
\]
\end{tcolorbox}
}

\vspace{1em}

\begin{examplebox}{เฉลยละเอียด}
มวลบนโต๊ะคือ

\[
m_1=3.0\,\si{kg}
\]

มวลแขวนคือ

\[
m_2=2.0\,\si{kg}
\]

พื้นโต๊ะลื่น รอกเบา และไม่มีแรงเสียดทาน

แรงภายนอกที่ทำให้ระบบเคลื่อนที่คือ น้ำหนักของมวลแขวน

\[
m_2g
\]

พิจารณาทั้งระบบ

\[
m_2g=(m_1+m_2)a
\]

ดังนั้น

\[
a=\frac{m_2g}{m_1+m_2}
\]

แทนค่า

\[
a=\frac{2g}{3+2}
\]

\[
a=\frac{2g}{5}
\]

มวลแขวนอยู่สูงจากพื้น

\[
s=2.0\,\si{m}
\]

มวลแขวนจะเร่งลงจนกระทบพื้น หลังจากนั้นเชือกไม่ดึงเพิ่ม ความเร็วสูงสุดของมวล $3.0\,\si{kg}$ จึงเกิดก่อนหรือขณะมวลแขวนกระทบพื้น

เริ่มจากหยุดนิ่ง

\[
v^2=2as
\]

แทนค่า

\[
v^2=2\left(\frac{2g}{5}\right)(2)
\]

\[
v^2=\frac{8g}{5}
\]

ดังนั้น

\[
\boxed{v=\sqrt{\frac{8g}{5}}}
\]
\end{examplebox}

% =========================================================
\subsection{ปี 2561 ตอนที่ 2 ข้อ 4: สองมวลบนพื้นเอียงลื่นชนกันกลางทาง}

\vspace{1.5em}

\IfFileExists{img/posn61s2q4.png}{
\begin{center}
    \includegraphics[width=1\linewidth]{img/posn61s2q4.png}
\end{center}
}{
\begin{tcolorbox}[title=ตำแหน่งภาพที่แนะนำ]
ให้ตัดภาพข้อสอบปี 2561 ตอนที่ 2 ข้อ 4 แล้วบันทึกเป็น

\[
\texttt{img/posn61s2q4.png}
\]
\end{tcolorbox}
}

\vspace{1em}

\begin{examplebox}{เฉลยละเอียด}
พื้นเอียงลื่นยาว

\[
d
\]

ทำมุม

\[
\theta
\]

กับแนวระดับ

มวล $A$ ถูกปล่อยจากหยุดนิ่งจากปลายบน และมวล $B$ ถูกดีดขึ้นจากปลายล่าง ทั้งสองชนกันที่กึ่งกลางพื้นเอียง

บนพื้นเอียงลื่น ความเร่งตามพื้นเอียงมีขนาด

\[
a=g\sin\theta
\]

\textbf{การเคลื่อนที่ของมวล $A$}

มวล $A$ เริ่มจากหยุดนิ่ง และเคลื่อนที่ลงพื้นเอียงเป็นระยะ

\[
\frac{d}{2}
\]

ใช้

\[
s=ut+\frac{1}{2}at^2
\]

โดย

\[
u=0
\]

ดังนั้น

\[
\frac{d}{2}=\frac{1}{2}(g\sin\theta)t^2
\]

จึงได้

\[
t^2=\frac{d}{g\sin\theta}
\]

\[
t=\sqrt{\frac{d}{g\sin\theta}}
\]

\textbf{การเคลื่อนที่ของมวล $B$}

ให้ความเร็วต้นของ $B$ ขึ้นตามพื้นเอียงเป็น

\[
u_B
\]

มวล $B$ เคลื่อนที่ขึ้นพื้นเอียงเป็นระยะ

\[
\frac{d}{2}
\]

แต่มีความเร่งลงตามพื้นเอียง ดังนั้น

\[
\frac{d}{2}=u_Bt-\frac{1}{2}(g\sin\theta)t^2
\]

จากสมการของมวล $A$ มี

\[
\frac{1}{2}(g\sin\theta)t^2=\frac{d}{2}
\]

แทนลงไป

\[
\frac{d}{2}=u_Bt-\frac{d}{2}
\]

ดังนั้น

\[
u_Bt=d
\]

\[
u_B=\frac{d}{t}
\]

แทนค่า $t$

\[
u_B=
\frac{d}{\sqrt{\frac{d}{g\sin\theta}}}
\]

\[
u_B=\sqrt{dg\sin\theta}
\]

ดังนั้น

\[
\boxed{u_B=\sqrt{dg\sin\theta}}
\]
\end{examplebox}

% =========================================================
\subsection{ปี 2565 ข้อ 18: ดึงกล่องบนพื้นหยาบด้วยแรงทำมุมลง}

\vspace{1.5em}

\IfFileExists{img/posn65q18.png}{
\begin{center}
    \includegraphics[width=1\linewidth]{img/posn65q18.png}
\end{center}
}{
\begin{tcolorbox}[title=ตำแหน่งภาพที่แนะนำ]
ให้ตัดภาพข้อสอบปี 2565 ข้อ 18 แล้วบันทึกเป็น

\[
\texttt{img/posn65q18.png}
\]
\end{tcolorbox}
}

\vspace{1em}

\begin{examplebox}{เฉลยละเอียด}
กล่องมวล $M$ อยู่บนพื้นหยาบ มีสัมประสิทธิ์เสียดทาน $\mu$ และถูกดึงด้วยแรง $F$ ทำมุม $\theta$ ใต้แนวระดับไปทางขวา

ต้องการหามุมมากที่สุดที่ยังสามารถลากกล่องไปทางขวาได้

แตกแรง $F$

องค์ประกอบแนวราบคือ

\[
F\cos\theta
\]

องค์ประกอบแนวดิ่งลงคือ

\[
F\sin\theta
\]

เพราะแรงมีองค์ประกอบกดลง แรงปกติจึงเป็น

\[
N=Mg+F\sin\theta
\]

แรงเสียดทานสูงสุดคือ

\[
f_{\max}=\mu N
\]

ดังนั้น

\[
f_{\max}=\mu(Mg+F\sin\theta)
\]

เพื่อให้ลากไปทางขวาได้ ต้องมี

\[
F\cos\theta>\mu(Mg+F\sin\theta)
\]

ย้ายพจน์ที่มี $F$

\[
F\cos\theta-\mu F\sin\theta>\mu Mg
\]

\[
F(\cos\theta-\mu\sin\theta)>\mu Mg
\]

ถ้าต้องการให้มีแรง $F$ บางค่าที่ทำให้ลากได้ ต้องมี

\[
\cos\theta-\mu\sin\theta>0
\]

ดังนั้น

\[
\cos\theta>\mu\sin\theta
\]

หารด้วย $\cos\theta$

\[
1>\mu\tan\theta
\]

\[
\tan\theta<\frac{1}{\mu}
\]

ดังนั้นมุมมากที่สุดคือ

\[
\boxed{\theta_{\max}=\tan^{-1}\left(\frac{1}{\mu}\right)}
\]
\end{examplebox}

% =========================================================
\subsection{ปี 2566 ข้อ 1: วงแหวนตกตามเชือกดิ่ง}

\vspace{1.5em}

\IfFileExists{img/posn66q1.png}{
\begin{center}
    \includegraphics[width=1\linewidth]{img/posn66q1.png}
\end{center}
}{
\begin{tcolorbox}[title=ตำแหน่งภาพที่แนะนำ]
ให้ตัดภาพข้อสอบปี 2566 ข้อ 1 แล้วบันทึกเป็น

\[
\texttt{img/posn66q1.png}
\]
\end{tcolorbox}
}

\vspace{1em}

\begin{examplebox}{เฉลยละเอียด}
วงแหวนมวล $m$ สวมรอบเชือกดิ่งและกำลังตกลงด้วยความเร่ง

\[
a
\]

ต้องการหาแรงเสียดทานที่เชือกกระทำต่อวงแหวน

แรงที่กระทำต่อวงแหวนมีสองแรงสำคัญ

\[
mg \text{ ลง}
\]

และ

\[
f \text{ ขึ้น}
\]

เพราะวงแหวนเคลื่อนที่ลง แรงเสียดทานจึงต้านการเคลื่อนที่และชี้ขึ้น

เลือกทิศลงเป็นบวก

\[
\sum F=ma
\]

ดังนั้น

\[
mg-f=ma
\]

จัดรูป

\[
f=mg-ma
\]

\[
f=m(g-a)
\]

ดังนั้น

\[
\boxed{f=m(g-a)}
\]
\end{examplebox}

% =========================================================
\subsection{ปี 2567 ข้อ 1: ไถลบนระนาบเอียงลื่นเทียบกับตกอิสระ}

\vspace{1.5em}

\IfFileExists{img/posn67q1.png}{
\begin{center}
    \includegraphics[width=1\linewidth]{img/posn67q1.png}
\end{center}
}{
\begin{tcolorbox}[title=ตำแหน่งภาพที่แนะนำ]
ให้ตัดภาพข้อสอบปี 2567 ข้อ 1 แล้วบันทึกเป็น

\[
\texttt{img/posn67q1.png}
\]
\end{tcolorbox}
}

\vspace{1em}

\begin{examplebox}{เฉลยละเอียด}
มวล $m$ ไถลจากหยุดนิ่งจากจุด $A$ บนระนาบเอียงลื่น $AB$ ถึงจุด $B$ ต้องหาเวลานี้เป็นกี่เท่าของการตกอิสระจาก $A$ ถึง $C$

จากรูป $A$, $B$, และ $C$ อยู่บนวงกลม โดย $AC$ เป็นรัศมีแนวดิ่ง และ $CB$ เป็นรัศมีแนวระดับ

ให้รัศมีวงกลมเป็น

\[
R
\]

ดังนั้น

\[
AC=R
\]

และ

\[
AB=R\sqrt{2}
\]

พื้นเอียง $AB$ ทำมุม

\[
45^\circ
\]

กับแนวระดับ

บนพื้นเอียงลื่น ความเร่งตามพื้นเอียงคือ

\[
a=g\sin45^\circ
\]

ดังนั้น

\[
a=\frac{g}{\sqrt{2}}
\]

วัตถุเริ่มจากหยุดนิ่งและเคลื่อนที่ตามพื้นเอียงเป็นระยะ

\[
s=R\sqrt{2}
\]

ใช้

\[
s=\frac{1}{2}at^2
\]

แทนค่า

\[
R\sqrt{2}
=
\frac{1}{2}\left(\frac{g}{\sqrt{2}}\right)t_{AB}^2
\]

จัดรูป

\[
t_{AB}^2=\frac{4R}{g}
\]

ดังนั้น

\[
t_{AB}=2\sqrt{\frac{R}{g}}
\]

ต่อไปหาเวลาตกอิสระจาก $A$ ถึง $C$

ระยะตกคือ

\[
R
\]

เริ่มจากหยุดนิ่ง

\[
R=\frac{1}{2}gt_C^2
\]

ดังนั้น

\[
t_C=\sqrt{\frac{2R}{g}}
\]

อัตราส่วนเวลาคือ

\[
\frac{t_{AB}}{t_C}
=
\frac{2\sqrt{\frac{R}{g}}}{\sqrt{\frac{2R}{g}}}
\]

\[
\frac{t_{AB}}{t_C}=\sqrt{2}
\]

ดังนั้น

\[
\boxed{\sqrt{2}}
\]
\end{examplebox}

% =========================================================
\subsection{ปี 2561 ตอนที่ 2 ข้อ 2: ก้อนกลมในระบบที่เร่งและผนังแนวตั้ง}

\vspace{1.5em}

\IfFileExists{img/posn61s2q2.png}{
\begin{center}
    \includegraphics[width=1\linewidth]{img/posn61s2q2.png}
\end{center}
}{
\begin{tcolorbox}[title=ตำแหน่งภาพที่แนะนำ]
ให้ตัดภาพข้อสอบปี 2561 ตอนที่ 2 ข้อ 2 แล้วบันทึกเป็น

\[
\texttt{img/posn61s2q2.png}
\]
\end{tcolorbox}
}

\vspace{1em}

\begin{examplebox}{เฉลยละเอียด}
โจทย์ให้ระบบเคลื่อนที่ด้วยความเร่ง

\[
a
\]

ไปทางซ้าย และทุกผิวไม่มีแรงเสียดทาน ต้องการหาแรง

\[
R
\]

ที่ผนังแนวตั้งกระทำต่อก้อนกลมมวล $m$

เลือกแกน $+x$ ไปทางขวา และแกน $+y$ ขึ้นด้านบน

เพราะก้อนกลมเคลื่อนที่ไปพร้อมกับระบบ จึงมีความเร่งในแนวราบไปทางซ้าย

\[
a_x=-a
\]

และไม่มีความเร่งในแนวดิ่ง

\[
a_y=0
\]

แรงที่กระทำต่อก้อนกลมมี 3 แรง

\[
mg \text{ ลง}
\]

\[
R \text{ จากผนังแนวตั้ง ชี้ไปทางซ้าย}
\]

และแรงปกติจากพื้นเอียง

\[
N
\]

พื้นเอียงทำมุม $\theta$ กับแนวระดับ ดังนั้นแรงปกติ $N$ ตั้งฉากกับพื้นเอียง และมีองค์ประกอบ

\[
N\sin\theta \text{ ไปทางขวา}
\]

\[
N\cos\theta \text{ ขึ้นด้านบน}
\]

พิจารณาแกน $y$

\[
\sum F_y=ma_y
\]

\[
N\cos\theta-mg=0
\]

ดังนั้น

\[
N=\frac{mg}{\cos\theta}
\]

พิจารณาแกน $x$

\[
\sum F_x=ma_x
\]

แรงทางขวาคือ $N\sin\theta$ และแรงทางซ้ายคือ $R$

\[
N\sin\theta-R=m(-a)
\]

\[
N\sin\theta-R=-ma
\]

จัดรูป

\[
R=N\sin\theta+ma
\]

แทนค่า $N=\frac{mg}{\cos\theta}$

\[
R=\frac{mg}{\cos\theta}\sin\theta+ma
\]

\[
R=mg\tan\theta+ma
\]

ดังนั้น

\[
\boxed{R=m(a+g\tan\theta)}
\]
\end{examplebox}

% =========================================================
\subsection{ปี 2565 ข้อ 5: ทรงกระบอกผิวลื่นบนแถบเชือก}

\vspace{1.5em}

\IfFileExists{img/posn65q5.png}{
\begin{center}
    \includegraphics[width=1\linewidth]{img/posn65q5.png}
\end{center}
}{
\begin{tcolorbox}[title=ตำแหน่งภาพที่แนะนำ]
ให้ตัดภาพข้อสอบปี 2565 ข้อ 5 แล้วบันทึกเป็น

\[
\texttt{img/posn65q5.png}
\]
\end{tcolorbox}
}

\vspace{1em}

\begin{examplebox}{เฉลยละเอียด}
ทรงกระบอกผิวลื่นมวล $M$ อยู่ในสมดุลบนแถบเชือกที่แขวนจากจุดต่างระดับดังรูป ต้องการหาแรงตึงในแถบเชือก

ให้แรงตึงเชือกแต่ละข้างมีขนาด

\[
T
\]

เนื่องจากระบบสมดุล แรงลัพธ์ต้องเป็นศูนย์

\[
\sum \vec{F}=0
\]

แรงที่กระทำต่อทรงกระบอกมีน้ำหนัก

\[
Mg
\]

ชี้ลง และแรงตึงจากแถบเชือกสองข้าง

แต่ละแรงทำมุม $\theta$ กับแนวระดับ ดังนั้นองค์ประกอบแนวดิ่งของแรงตึงแต่ละข้างคือ

\[
T\sin\theta
\]

องค์ประกอบแนวราบของแรงตึงสองข้างหักล้างกัน เพราะมีขนาดเท่ากันและทิศตรงข้ามกัน

พิจารณาแกนดิ่ง

\[
\sum F_y=0
\]

แรงขึ้นจากเชือกสองข้างรวมกันเป็น

\[
2T\sin\theta
\]

จึงได้

\[
2T\sin\theta-Mg=0
\]

ดังนั้น

\[
2T\sin\theta=Mg
\]

\[
T=\frac{Mg}{2\sin\theta}
\]

จึงได้

\[
\boxed{T=\frac{Mg}{2\sin\theta}}
\]
\end{examplebox}

% =========================================================
\subsection{ปี 2565 ข้อ 8: โซ่สมดุลและแรงในแนวระดับ}

\vspace{1.5em}

\IfFileExists{img/posn65q8.png}{
\begin{center}
    \includegraphics[width=1\linewidth]{img/posn65q8.png}
\end{center}
}{
\begin{tcolorbox}[title=ตำแหน่งภาพที่แนะนำ]
ให้ตัดภาพข้อสอบปี 2565 ข้อ 8 แล้วบันทึกเป็น

\[
\texttt{img/posn65q8.png}
\]
\end{tcolorbox}
}

\vspace{1em}

\begin{examplebox}{เฉลยละเอียด}
โซ่หรือเชือกมวล $M$ แขวนอยู่ในสมดุลระหว่างจุด $O$ กับ $A$ ต้องการหาแรงในแนวระดับ

\[
H
\]

จากรูป ที่ปลาย $A$ แรงตึงเชือกมีทิศตามแนวเส้นสัมผัส และทำมุม

\[
\theta
\]

กับแนวระดับ

ให้แรงตึงที่ปลาย $A$ มีขนาด

\[
T
\]

พิจารณาโซ่ทั้งเส้นเป็นระบบเดียว แรงภายนอกที่กระทำต่อระบบมี 3 แรง

\[
Mg \text{ ลง}
\]

\[
H \text{ ตามแนวระดับ}
\]

และ

\[
T \text{ ตามแนวสัมผัสที่จุด } A
\]

เพราะโซ่อยู่ในสมดุล

\[
\sum F_x=0
\]

\[
\sum F_y=0
\]

แตกแรงตึง $T$ เป็นองค์ประกอบ

\[
T_x=T\cos\theta
\]

\[
T_y=T\sin\theta
\]

พิจารณาแกนดิ่ง

\[
T\sin\theta-Mg=0
\]

ดังนั้น

\[
T=\frac{Mg}{\sin\theta}
\]

พิจารณาแกนนอน

\[
H=T\cos\theta
\]

แทนค่า $T$

\[
H=\frac{Mg}{\sin\theta}\cos\theta
\]

\[
H=Mg\cot\theta
\]

ดังนั้น

\[
\boxed{H=Mg\cot\theta}
\]
\end{examplebox}
